% !TeX root = ../papers/nc-p3.tex
\section{五中心商の乗法表から全ての四中心商を求める}
\label{sec:rank-fifteen-center-border}
\begin{proposition}\label{prop:five-center-border}
四周期の二次経路代数 $A$ と五つの中心四次式 $z_1,\ldots,z_5$ を取り、
固定した始点で $M=A/(z_1,\ldots,z_5)$ の長さ8が零とする。
$[b_1:\cdots:b_5]\in\mathbb P^4$ に垂直な四式で割った始点加群を $Q_b$ とする。
$M$ の長さ7以下の語基底 $W$ を用い、次の一つの多項式行列の計算を全ての $b$ に使える。
\begin{enumerate}
\item $F=\bigoplus_{w\in W}k[Z]w$, $\deg Z=4$ と置く。
長さ8までの語消去で各矢の右作用を $F$ 上に持ち上げる。
除算するのは五中心商の語消去の主元だけで、$b$ による除算はない。
\item 持ち上げた二次関係を全基底語に作用させ、その残余と、最初の五中心関係の残余を取る。
$Z$ と全矢の右作用でそれらの線型包を閉じたものを $N$ とすれば $Q_b=F/N$ である。
従って各次数への伝播は、長さ7以下の語に対する固定した乗法表だけを使う。
\item 五中心商の各文字の次数0から7の次元が
\[
 (1,0,0,0),\ (1,1,1,1),\ (2,2,2,4),\ (5,5,5,5),\
 (6,8,8,8),\ (9,9,9,9),\ (10,10,10,4),\ (5,5,5,5)
\]
とする。四中心商の長さ9から13の文字別の目標上界を順に
$(5,5,5,5),(2,2,2,4),(1,1,1,1),(1,0,0,0),(0,0,0,0)$
とすると、各条件は $b$ の斉次小行列式の非消滅の和であり、その次数上界は順に
\[
 (6,6,6,6),\ (12,12,12,8),\ (14,14,14,14),\ (13,16,16,16),\
 (21,21,21,21)
\]
となる。ここでは五中心が文字0の場合を扱う。
\end{enumerate}
\end{proposition}
\begin{proof}
$Q_b$ は中心変数 $Z$ を加え $z_j=b_j Z$ を課した代数の始点行である。
実際、ある $b_i\ne0$ なので $Z$ を $z_i/b_i$ で消去できる。
$Z$ を割った商が $M$ であり、次数帰納で $F\to Q_b$ が全射である。

乗法表を具体化する。長さ $n$ でまだ決まらない列は、長さ $n-1$ の五中心商の
基底語に一文字を付けたものだけである。その他は $Z F_{n-4}$ に入り、既出の表を使う。
長さ $n-2$ の基底語に末尾二次関係を付けた行、および $n=4$ では
五行 $z_j-b_j Z$ を作る。新しい列だけへの射影は $M$ の語の関係行列そのものである。
これは $Z=0$ とすれば、各段で前段までの関係と最後の二次関係を全て課した商となることによる。
中央関係の右倍は前段の乗法の簡約に含まれ、中央性から途中へ挿入した関係も含まれる。
主元列を非主元列と $Z F_{n-4}$ の和で表し、残る列を長さ $n$ の基底語とする。
$M_8=0$ なら長さ8では全ての新しい列が消える。$Z$ と可換に延ばして全基底語の
右作用が定まった。全ての式は実際の $Q_b$ で成り立つ。

新しい列を消去した後の残余も関係として保持する。$F/N$ 上では全基底語で
二次関係が零なので、矢の作用は $A$ の右作用を与える。空語に矢を掛ければ全非主元語を回収できる。
$Z$ は作用と可換で、$z_j$ は $A$ の中心だから、空語で確認した
$z_j=b_j Z$ は全ての $Z^d w$ にも成立する。従って $Q_b\to F/N$ が得られ、
基底とその中心倍で両者の自然な写像は互いに逆である。
長さ8までに残余がなく、上の次元の場合には、最初の関係は長さ7の基底語に
二次関係を掛けた九次の残余だけで尽きる。これ以後はこれらの右倍と $Z$ 倍を取ればよい。

長さ $n$ の一行で、列 $Z^j w$ の係数は $b$ について次数 $j$ の斉次式となる。
実際、中央関係の右辺で $Z$ と一緒に $b_i$ が一つ加わり、元の二次関係は $b$ を含まない。
既出の行の $Z$ 倍は行全体で次数が一つ下がる。同次な全右倍をそのまま並べるか、
$Z$ 倍の代わりに全ての $b_iZ$ 倍を並べれば射影点での線型包は変わらない。
必要な階数を $r$ とすれば、$r$ 列の $j$ の和の最大値が小行列式次数の上界になる。
例えば十三次の各文字は $Z^2 M_5$ に九列、$Z^3 M_1$ に一列あり、
全階数の小行列式は $2\cdot9+3=21$ 次以下である。各次数で同じ計算を行うと表示した表を得る。
従って各文字・次数である良い $b$ が存在するかの検査には、四変数・全次数21の
補間格子で足りる。アフィン座標に相異なる22標本を用いるNewton補間により、
正標数でも必要なら有限体拡大で同じ結論が従う。
\end{proof}

\begin{proposition}\label{prop:ordinary-bivariate-border-cut}
標数3の非分岐単元表示の関数体 $\mathbb F_3(K,H)$ で、始点0において
五本の規格化中心の第2,3,4,5列による商は、長さ1から13まで
\[
               4,10,20,31,40,44,40,31,20,10,4,1,0
\]
の次元を持つ。この証明書の除外因子は元の $(K,H)$ 平面上で171個である。
\end{proposition}
\begin{proof}
前命題の二次関係は語の偶奇による二次拡大なしに持ち上げられる。
$J$ を六平方の積、$\alpha^2=J$ とし、次の矢 $a$ について
\[
 \delta(n,h,a)=\bigl([p+n=2\bmod4,\ a=3]+\beta(h,a)\bigr)\bmod2
\]
と置く。これは文字 $h$ の長さ $n$ の語に $a$ を付けたときの
$\ell_p$ の増分の偶奇であり、$n$ を4減らしても同じである。
そこで $a$ の右作用を $\alpha^\delta$ 倍して用いる。この正規化で
語 $w$ の係数の基底変換は $\alpha^{D(w)}$、ただし $D(w)$ は順次の $\delta$ の和である。
中心式で必要な追加倍率 $\alpha^{\ell_p(w)-D(w)}$ の指数は偶数である。
二次関係の各係数の $\alpha$ の指数から連続二回の $\delta$ を引いても、同一行では同じ偶奇となる。
従って行全体の共通倍率を除き、全ての新旧列を $\mathbb F_3(K,H)$ 上で扱える。

$Z$ 倍で短い語に戻る列を保存したまま、五中心商の新しい語の列だけを掃き出す。
初期の長さ8まで余分な関係はなく、九次には各文字の自由十列に独立な五つの残余が現れる。
以後は直前の残余の原行を選んで一文字付け、四つ前の残余の中心倍を併合する。
九次以後の文字別商は前命題の目標値にちょうど一致する。
\path{computations/verify_rank_fifteen_bivariate_border.py} は二変数のままこれを検算する。
主元の分子、最初の経路と中心の単元を含めて171既約因子が現れる。
各因子の零点を除けばこの始点での上界は特殊化できる。
さらに \path{computations/verify_rank_fifteen_center_border.py} は
前出の $\mathbb F_{81}$ の点の七つの中心方向と四始点で、全長について
元の四中心商から始める直接消去と照合する。
この関数体計算に残る例外点での正則性は、関数体の次元からは決まらない。
\end{proof}
